\documentclass[../main]{subfiles}
\begin{document}
\chapter{Ciclos termodinámicos}
\img{images/06/ciclos.jpg}{Sección Transversal pistones del motor. }
\info{
Contenido}{
Ciclos termodinámicos: Diésel,Otto,Brayton y Rankine.
}

\info{Ciclo termodinámico}{
Se denomina ciclo termodinámico a cualquier serie de procesos termodinámicos tales que, al transcurso de todos ellos, el sistema regresa a su estado inicial; es decir, que la variación de las magnitudes termodinámicas propias del sistema se anula.
}
\actividad{Trabajo de investigación}
\begin{enumerate}
    \item Investigar y explicar el ciclo diésel.
    \begin{enumerate}
        \item Biografía de Rudolf Diesel.
        \item Funcionamiento del ciclo Diesel:admisión, compresión, combustión y escape. Graficar.
    \end{enumerate}
     \item Ciclo Naftero o ciclo Otto.
    \begin{enumerate}
        \item Explicar el ciclo Otto. Graficar.
        \item explicar las fases: Admisión a presión constante, compresión isoentrópica, combustión y trabajo, escape.
    \end{enumerate}
    \item Explicar el ciclo de Brayton.
    \begin{enumerate}
        \item Escribir la biografía de George Bailey Brayton.
        \item Explicar los tipos de ciclo de Brayton: Abierto, cerrado y inverso. Graficar cada caso.
    \end{enumerate}
    \item Ciclo de Rankine.
    \begin{enumerate}
        \item Explicar el ciclo de Rankine.
        \item Tipos de ciclo de Rankine: Abierto, cerrado, explicar partes del proceso y graficar diagrama de proceso(P-V).
    \end{enumerate}
\end{enumerate}



\end{document}